\section*{Chapitre \ref{ch:b3:trace-element-modelling} --- Modélisation des éléments en traces}

\begin{solution}{exo:b3:trace-element-modelling:1}
Lherzolite à spinelle~: $0.12\times0.40 + 0.88\times0.016 = 0.048 + 0.0141 =
0.062$. Lherzolite à grenat~: $0.195\times0.40 + 0.167\times0.11 +
0.639\times0.016 = 0.078 + 0.0184 + 0.0102 = 0.107$.
\end{solution}

\begin{solution}{exo:b3:trace-element-modelling:2}
$1/(0.05 + 0.02\times0.95) = 1/0.069 = \qty{14.5}{ppm}$;
$1/(0.05 + 0.10\times0.95) = 1/0.145 = \qty{6.9}{ppm}$.
\end{solution}

\begin{solution}{exo:b3:trace-element-modelling:3}
L'enrichissement est maximal lorsque $F \to 0$~: $1/D = 100$. Avec $D = 0$,
$C_L/C_0 = 1/F = 20$.
\end{solution}

\begin{solution}{exo:b3:trace-element-modelling:4}
$\mathrm{La/Yb} = 30/2.5 = 12.0$; $(\mathrm{La}/\mathrm{Yb})_N =
(30/0.232)/(2.5/0.163) = 129.3/15.34 = 8.4$.
\end{solution}

\begin{solution}{exo:b3:trace-element-modelling:5}
\omterm{thm:b3:trace-element-modelling:batch}{Fusion à l'équilibre}~: $C_S/C_0 = 0.01/(0.01 + 0.05\times0.99) = 0.168$. \omterm{thm:b3:trace-element-modelling:fractional}{Fusion fractionnée}~:
$(0.95)^{99} = 0.0062$, soit 27 fois moins. Le liquide instantané contient
$0.0062/0.01 = 0.62\,C_0$~: moins que la source.
\end{solution}

\begin{solution}{exo:b3:trace-element-modelling:6}
$D_0 = 0.58\times0.0003 + 0.27\times0.002 + 0.12\times0.053 + 0.03\times0.01
= 0.00737$; $P = 0.20\times0.0003 + 0.25\times0.002 + 0.45\times0.053 +
0.10\times0.01 = 0.0254$. Pour $F = 0.02$~: \omterm{def:b3:trace-element-modelling:nonmodal}{fusion non modale} $1/(0.00737 +
0.02\times0.9746) = 37.2$, fusion modale $1/(0.00737 + 0.02\times0.9926) = 36.7$. Pour
$F = 0.10$~: 9.54 et 9.38, soit un écart de \qty{1.7}{\percent}.
\end{solution}

\begin{solution}{exo:b3:trace-element-modelling:7}
$D = 0.1$~: \omterm{thm:b3:trace-element-modelling:batch}{fusion à l'équilibre} $1/(0.1 + 0.09) = 5.26$~; \omterm{def:b3:trace-element-modelling:fractional}{liquide de fusion fractionnée agrégé} $(1 - 0.9^{10})/0.1 =
6.51$. $D = 0.01$~: \omterm{thm:b3:trace-element-modelling:batch}{fusion à l'équilibre} $1/(0.01 + 0.099) = 9.17$~; \omterm{def:b3:trace-element-modelling:fractional}{liquide de fusion fractionnée agrégé}
$(1 - 0.9^{100})/0.1 = 10.00$. Le couple à $D = 0.1$ diffère davantage~: pour
$D = 0.01$, les deux liquides contiennent déjà presque tout l'élément, et tendent tous deux vers
$1/F$.
\end{solution}

\begin{solution}{exo:b3:trace-element-modelling:8}
$D' = 0.99\times0.001 + 0.01 = 0.01099$. Premier liquide~: $1/D' = 91$, contre
1000 en \omterm{thm:b3:trace-element-modelling:fractional}{fusion fractionnée}. Pour $F = 0.01$~: $91\times0.99^{90.0} = 36.8$,
contre $1000\times0.99^{999} = 0.044$~: le liquide interstitiel empêche le résidu
d'être dépouillé de ses éléments.
\end{solution}

\begin{solution}{exo:b3:trace-element-modelling:9}
Le rapport est multiplié par $(D_{\mathrm{Sm}} + F(1 - D_{\mathrm{Sm}}))/(D_{\mathrm{La}}
+ F(1 - D_{\mathrm{La}}))$~: pour $F = 0.01$, $0.0714/0.0177 = 4.0$~; pour $F = 0.10$,
$0.156/0.107 = 1.46$. Le rapport est sensible à $F$ lorsque $F$ est faible, et presque
insensible au-delà d'environ \qty{10}{\percent}~; il dépend aussi de $D_{\mathrm{Sm}}$,
donc de la nature des minéraux résiduels.
\end{solution}

\begin{solution}{exo:b3:trace-element-modelling:10}
$\frac1F\int_0^F \frac1D(1 - x)^{1/D - 1}\,\mathrm dx = \frac1F\bigl[-(1 -
x)^{1/D}\bigr]_0^F = [1 - (1 - F)^{1/D}]/F$, la formule du chapitre~; si
$(1 - F)^{1/D} \to 0$, elle tend vers $1/F$ fois $C_0$. À moins de \qty{1}{\percent} près~:
$(1 - F)^{100} < 0.01$, soit $F > 1 - 0.01^{0.01} = 0.045$, environ
\qty{4.5}{\percent}.
\end{solution}

\begin{solution}{exo:b3:trace-element-modelling:11}
Pente $a = (10 - 6)/(20 - 10) = \qty{0.40}{ppm^{-1}}$~; ordonnée à l'origine $b = 10 -
0.40\times20 = 2.0$. $D_{\mathrm{Yb}} = 0.40\times0.44 = 0.176$;
$C_0^{\mathrm{La}} = 2.0\times0.44/0.824 = \qty{1.07}{ppm}$. $F_P = 1.07/20 =
\qty{5.3}{\percent}$, $F_Q = 1.07/10 = \qty{10.7}{\percent}$.
\end{solution}

\begin{solution}{exo:b3:trace-element-modelling:12}
$F = x_i/p_i$~: clinopyroxène $0.12/0.45 = 0.27$, spinelle $0.03/0.10 = 0.30$,
orthopyroxène $0.27/0.25 = 1.08$, olivine $0.58/0.20 = 2.9$~; clinopyroxène
fond en premier. Pour $F = 0.2$~: $\bar D = (0.00737 - 0.0254\times0.2)/0.8 =
0.0029$. Au-delà de $F = 0.27$, les $p_i$ changent (il ne reste plus de clinopyroxène à fondre),
et la formule, qui suppose des $p_i$ constants, ne s'applique plus.
\end{solution}

\begin{solution}{pb:b3:trace-element-modelling:1}
\textbf{1.} $D_{\mathrm{La}} = 0.12\times0.058 + 0.88\times0.001 = 0.00784$;
$D_{\mathrm{Yb}} = 0.12\times0.58 + 0.88\times0.16 = 0.210$.

\textbf{2.} $D_{\mathrm{La}} = 0.195\times0.058 + 0.167\times0.001 +
0.639\times0.001 = 0.0121$; $D_{\mathrm{Yb}} = 0.195\times0.58 +
0.167\times4.5 + 0.639\times0.16 = 0.967$. Valeurs arrondies~: 0.0078, 0.21, 0.012, 0.97.

\textbf{3.} $D_{\mathrm{Yb}}$, de 0.21 à 0.97, parce que le grenat retient les
terres rares lourdes ($D = 4.5$). L'ytterbium reste encore très faiblement incompatible,
$0.97 < 1$.

\textbf{4.} $(0.648/0.232)/(0.441/0.163) = 2.79/2.71 = 1.03$~: un diagramme normalisé plat
environ 2.8 fois la chondrite.

\textbf{5.} $0.648/(0.0078 + 0.01\times0.9922) = \qty{36.6}{ppm}$~; à
\qty{5}{\percent}, \qty{11.3}{ppm}~; à \qty{10}{\percent}, \qty{6.05}{ppm}.

\textbf{6.} $F = (0.648/9.6 - 0.0078)/0.9922 = (0.0675 - 0.0078)/0.9922 =
0.0602$, \qty{6.0}{\percent}.

\textbf{7.} $\mathrm{Yb} = 0.441/(0.21 + 0.0602\times0.79) = \qty{1.71}{ppm}$
et La/Yb $= 5.6$~; la lave A donne 1.7 et 5.6~: elle correspond à une source à spinelle.

\textbf{8.} La $= 0.0078\times9.6 = \qty{0.075}{ppm}$, Yb $= 0.21\times1.71 =
\qty{0.36}{ppm}$. Il reste du lanthane~: $0.940\times0.075/0.648 = 0.11$, soit environ
\qty{11}{\percent}.

\textbf{9.} $F = 0.648/9.6 = \qty{6.75}{\percent}$.

\textbf{10.} $(0.648/0.0078)(1 - F)^{127.2}$~: \qty{73}{ppm} à
\qty{0.1}{\percent}, \qty{23}{ppm} à \qty{1}{\percent}, \qty{0.03}{ppm} à
\qty{6}{\percent}.

\textbf{11.} $0.648\times(1 - 0.940^{128.2})/0.060 = \qty{10.8}{ppm}$.

\textbf{12.} $F = \qty{6.75}{\percent}$~: à ce $F$, $(1 - F)^{1/D} =
\num{1.3e-4}$, donc le \omterm{def:b3:trace-element-modelling:fractional}{liquide de fusion fractionnée agrégé} contient tout le lanthane de la source, et
$\bar C_L = C_0/F$ comme pour un élément parfaitement incompatible.

\textbf{13.} $0.648\times0.9325^{127.2} = \qty{9e-5}{ppm}$, soit environ 800 fois
moins que les \qty{0.075}{ppm} de la \omterm{thm:b3:trace-element-modelling:batch}{fusion à l'équilibre}.

\textbf{14.} Les deux modèles encadrent la réponse~: \qty{6.0}{\percent} (\omterm{thm:b3:trace-element-modelling:batch}{fusion à l'équilibre})
à \qty{6.75}{\percent} (\omterm{thm:b3:trace-element-modelling:fractional}{fusion fractionnée} agrégée). Le lanthane fixe $F$ à
environ \qty{1}{\percent} près, quel que soit le modèle.

\textbf{15.} A~: La/Yb $= 5.6$, $(\mathrm{La}/\mathrm{Yb})_N =
(9.6/0.232)/(1.7/0.163) = 4.0$. B~: La/Yb $= 34.7$, $(\mathrm{La}/\mathrm{Yb})_N
= (15.6/0.232)/(0.45/0.163) = 24.4$.

\textbf{16.} Lorsque $F \to 0$~: $(0.648/0.0078)/(0.441/0.21) = 83.1/2.10 = 39.6$.

\textbf{17.} $F = \qty{0.11}{\percent}$, où le liquide contiendrait
\qty{72}{ppm} de La et \qty{2.1}{ppm} d'Yb. La lave B contient \qty{15.6}{ppm} de La~: une
source à spinelle ne peut la produire.

\textbf{18.} $F = (0.648/15.6 - 0.012)/0.988 = 0.0299$, \qty{3.0}{\percent}~;
$\mathrm{Yb} = 0.441/(0.97 + 0.0299\times0.03) = \qty{0.454}{ppm}$, comme
observé.

\textbf{19.} Le grenat remplace le spinelle vers 60 à \qty{90}{km}~: B a commencé à
fondre plus profondément que cela, A moins profondément, dans le domaine du spinelle. Une fusion de lherzolite à grenat
ayant \qty{9.6}{ppm} de La ($F \approx \qty{5.6}{\percent}$) aurait un rapport La/Yb
voisin de 21, et non 5.6~: A ne provient pas de la lherzolite à grenat.

\textbf{20.} $a = 0.97/0.441 = \qty{2.20}{ppm^{-1}}$, $b = 0.648\times0.03/0.441
= 0.044$~; pour La $= 15.6$, La/Yb $= 2.20\times15.6 + 0.044 = 34.4$, contre
34.7 pour la lave B.

\textbf{21.} L'olivine rejette les deux éléments~: le rapport La/Yb est inchangé. Le liquide primaire
avait $9.6\times0.9 = \qty{8.64}{ppm}$ de La, donc $F = (0.648/8.64 -
0.0078)/0.9922 = \qty{6.8}{\percent}$ au lieu de 6.0.

\textbf{22.} Une seule source peut produire les deux laves~: B est un liquide de faible taux de fusion
(\qty{3}{\percent}) qui a commencé à se former dans le domaine du grenat, plus profondément que vers 60 à
\qty{90}{km}~; A est un liquide de taux de fusion plus élevé, formé dans le domaine du spinelle.

\textbf{23.} La lave A résulte d'une
\omterm{thm:b3:trace-element-modelling:batch}{fusion à l'équilibre} à \textbf{\qty[text-series-to-math=true]{6.0}{\percent}} d'une lherzolite à spinelle (\qty{6.75}{\percent} si le liquide a été
extrait par \omterm{thm:b3:trace-element-modelling:fractional}{fusion fractionnée} puis agrégé).
\end{solution}
